
\subsection{Summary of Contributions}
Summarize main findings and improvements over baseline.

\subsection{Future Work}
Discuss architectural improvements. Discuss scalability.
Discuss transformers and potential benefits with hardware specifically made for this. 

Explore integration into high-level trigger.

Discuss training per epoch being a bit wierd. \textbf{Already done?}

Discuss trimming the model to minimal size but preserving performance, a study that finds the optimal models in this sense as well. This is not something that is of interest until very late. This is only interesting for production, and even then production may have high computational capabilities anyway. As phycisist we are most interested in accuracy first and foremost. We are interested in finding the limit of ambiguity: the boundary of maximum performance where we simply cannot push further than what the inherent ambiguity is in the event. If we find this boundary we can also use this info to know what poisson sampling parameter we can set the whole experiment. It gives us the information what the optimal level of pile-up is. As discussed in background, pile-up presents a classifcal trade-off engineering problem. If we have optimal models that squeeze out as much information as possible that is given to us, then we can adjust this trade-off finding an optimal level, thus optimizing the physics reach. The MLPF model was motivated by the incoming LHC high luminosity phase for example, where pile-up will increase to unprecidented levels in their experiment, then at this level ML was interpreted as a potentially valuable tool. 

Discuss how MORE information can be fed and how that has not been tested yet. Discuss that MORE information can mean new unique information but it can also include deriving more object-level information that is already processed beforehand somewhat but maybe the models are not able to construct themselves. 

This inlcudes: deriving explicit information related to incident momentum angles (this metric is basically totally absent for the scope of my thesis), exploring new object level information such as total measured ECal energy (even though it technically is available we do not know if the models pick up on this, this is due to the black-box nature of ML models), deriving angle information that is derived from across the geometrical space between the tpad and the ecal, implementing a fix in ldmx-sw that enables to extract origin ID labels for the tpad objects (centroids) as that is not available currently (then one could also make an across tpad-ecal hit assignment). One could also include the hcal subdetector exploring completely different assignments, such as a veto, signal event veto detection, or other the sky is the limit really. This is basically exploring FULL EVENT reconstruction, also similar to the idea of MLPF. 

Discuss the inclusion of TS and tpad information. Why heterogenous detector information is interesting, and discuss why the tpad information in this case seem to matter little to nothing at all for one task, and discuss why it might be logical that it matters more for the electron counting task and how the results slightly point in this direction, though it seems that the high granularity ecal information does go a long way in itself. Discuss what that could mean, discuss if we think that we can squeeze more infomration out or not. When do we expect this information to be important? When is it not important? 

One can explore deconstructing the black-box nature of the ML models finding ways to measure what information it finds valuable. I have only made a small test with the tpad ablations trying to use that to assess if the tpad is valuable for certain tasks, but one can make more fine grained analysis exploring what information from the ecal is most valuable and so on. There are ways to try to identify "what the models are actually doing". They are black-boxes and it exists an internal structure, they are not "magic-boxes" where there would be no internal structure just magic. 

Discuss the 3 highlighted simplifications and why they were made. Give some insight to the history. It was a very useful practical approach deconstructing the problem to perhaps its simplest possible way, then it allowed me to develop it further where it laid a path for the next step stepping stone introducing new tasks and taking inspiration for more sophisticated architectures. 

Discuss the comparison between transformer and gnn based models and why transformers may seem to be the preffered choice given this study. Transformers are winning the hardware lottery, there seem to be indication that the industry shift towards this archictecture is valid for HEP as well as many other industries including NLP related tasks. Transformers simply scale extremely well with hardware and thus train faster and perform inference faster. Sure they take up more storage but that is not the limiting factor the storage is so small so it is basically a constant negliable effect. This study showed that the baseline transformer model had an edge in accuracy (this is the defining metric after all, everything else does not matter, in offline analysis at least) and this was more present for the 3e event. That might indicate that transformers are able to utilize the global event context in some way that the gravnet models are restrictied to their local neighbourhoods. In that sense, then one could just increase hyperparameter k of gravnet, but scaling that up to fully connected graph would be counterintuitive and it would break the whole purpose and idea of using the gravnet in the first place then in this scenario the Transformer is just the supoerior architecture for fully connected graphs. The most important aspect in the end is the accuracy, the model that is able to deconstruct all the parts of the input collection and that performs the best reconstruction accuracy will be the model of interest, all else is below. This provided ml\_ldmx framework has built one example of an implementation that should allow for further studies with both the gravnet and the transformer. One could also explore trying to use the particlenet model, which has been a successful gravnet based model, but it was built for jet tagging tasks and it is unclear how well that translates to ldmx. 

It is important to understand why the transformer and gnn gravnet perform so similar, they are both set based neural networks and they essentially both belong to the gnn class of neural network models, they more or less approach the same type of problem in slightly different ways. 

Discuss the prospect of general optimization trimming the size of the network: parameters, inference time, size on disk. All these aspects probably only exclusively matter for online reconstruction or for a high level trigger task where time matters for meeting the deadline. 

Discuss prospect of adding tasks to the model (multi-task already or not). Are there new tasks that could be introduced that in post-processing can enhance the results even further? Other creative ideas? 

Discuss the most obvious and most interesting extensions that has been found in the thesis so far: 

electron counting: this is most relevant for triggering and for online reconstruction and this study has shown promise for future work here. If I had more time and a larger scope (this is merely 30 credits after all I REQUIRE to at least mention this once) maybe the first thing I would do is to create a dataset of 1e, 2e, ... 5e electrons maybe and try to approach a general-multiplicity electron counting model. I would also switch over to ecal trigger cells and see what the maximum performance would be here and see how far one can go with this really making a good effort of maximizing performance. This also has the potential to be the most rewarding and yielding result if it is good because it has a direct connection and positive result on the physics reach and "experimental efficiency" if it means that we can turn up the poisson sampling of the experiment and that we can allow for more pile-up! 

Discuss the canonical y-ordering. It started out from a very practical-need when we bumped into the problem that labels were swapped and had no consistent ordering in extracting root files. This jumbled the loss since label ordering matters for loss so the solution was to create a simple orderin technique that made things consistent. This "solution" with canonical y-ordering was maybe more advanced in the end because there was still a need to introduce permutation invariant accuracy, a quite obvious problem really. Herein also lies the solution really: hungarian matching, or similar, it was even clearly mentioned in one of my sources that collects how gnn:s can be used in HEP \cite{duarte2020-gnn-tracking}. Hungarian matching can be used to calculate loss, thus removing that the model is punsished if it make a good grouping but switched up labels, then only grouping matter invariant of label ordering. I did try to use this loss function instead but I did not preliminary get any visible improvement so it was abandoned. Sometimes models are able to brute-force their way through a problem that was unnecesarily hard for them, since the global minimum optimal solution is still to "just do things correct". This might be the underlying reason, so now the model gets an unnecessarily difficult task, but this task maybe still have the same optimal solution, even though one might think there is a risk that the model becomes to cautious in predicting if the labeling might be confused. The most physically interpretable is in the end the hungarian matching so that should definitely be explored in future work probably replacing canonical y ordering for loss, but the canonical y ordering is still useful for finding a consistent ordering it was the reason for why the confusion matrices hinted at that shower distance had a clear effect on the performance for example. 

The hungarian matching test that did not yield better results reminds me of an earlier implementation that I had for hard labeling dominant origin label. It was labeled by the maximum individual contribution, even when the contribs in one hit had several contribs of the same electron label! This mean that the summed energy could be larger from label 1, but label 2 had a single contribution larger than the other one in isolation. This does not make sense at all from a physics standpoint and it seems like an unnecesarily hard task for the model. However, when I realised this fault and corrected it it had no visible effect on the accuracy. It asks the question if we infer an induction bias as humans just to make things better interpretable... 

Discuss implementing LSH from mlpf 2021 that it used to make their gnn gravnet models in that report more efficient, or explore implementing flash attention as with mlpf latest iteration transformer 2026 for speed up and efficiecny (only relevant for online reconstruction probably). Discuss why this is probably overkill for this project with relatively few inputs. 

Discuss just investigating the mlpf model further, this work has been substantial for the scope of 30 credits and there simply was not so much time left in taking the next step and making this mlpf inspired multi task model, thus the "experimental" caption. Understanding and breaking mlpf down better might give new insights and improvements on my models. 

Put emphasis on that the codebase that was the basis for this report is meant to be logically structured and documented in such a way so that it can be handed over and that it can act as a platform for deep learning ML based tasks connected to LDMX. One of the intended purposes of the master thesis was to create something that could be picked up at a later stage. It should be a platform where resources can be reused to see if ML based methods can perform  different reconstruction tasks, by students or researchers alike. There had been many students before me that had done similar ML related work, but their codebase did not provide a solution to the first steps of reading the root files, extracting and tensorizing. I want to help the next student skip this part, as it is not as interesting from a research standpoint, and then more time could be spent on pushing the frontier for ldmx! 

Mention that there is an alternative geometry suggestion for trigger scintillator that introduces measurements in the x-direction and state that it can be used to achieve resolution in both x and y directions, thus yielding new unique information for the models that they might utilize.

